\textbf{Threats to validity.} \emph{Construct:} AP, HOTA and crossing F1 measure agreement with teacher labels;
human labels (protocol above) are the remedy, and the label-bias, GT-linking and label-audit analyses bound but do
not remove the problem. \emph{Internal:} the students differ in architecture, resolution and schedule, so only the
matched-resolution rows (Sec.~\ref{sec:res}) compare architectures; the adaptation variants share one recipe and one
seed, so differences under about 0.01 AP are not interpreted; the guardrails and the policy layer were designed
after seeing the E3 utterances, which are therefore reported separately from the E6 test cases written before any
E6 run; the same authors wrote the policy layer and the E6 cases, and no adaptive red-teaming against the final
policy was done, so E6 rejection rates are an upper bound on what a determined attacker would face. \emph{Statistical:} block-bootstrap intervals describe frame sampling within one video; with five 10\,s
blocks in the aa1 test split no test can be very small, and none of our intervals covers scene-to-scene variation.
\emph{External:} four scenes (one city plaza, three campus locations, three of them phone or handheld video) are
still few, and the three campus scenes appear to share one population; no night, rain or overhead camera is included; the
VLM results hold for one 2B model at temperature~0 and one prompt set.
